% ===== §14 Dormancy: The Inward Phase =====
\section{Dormancy: The Inward Phase}
\label{sec:rose}

\noindent
The spiral has a descending arc. There are stretches of a generative relation in which little appears to be generated---in which the partners have drawn inward, the throughput has slackened, the visible signs of coupling have gone quiet. Judged against the demand for unbroken concordance, such a phase reads as failure or as the beginning of an end. This short section argues that it may read instead as \emph{dormancy}: the low, inward phase of a driven system in which force is not lost but gathered, and from which the next ascent of the spiral is drawn. The section is the last of the cosmological movement, and it is also where the paper allows itself, briefly, to return to the image with which it began.

\subsection{The low-activity phase is not decline}
\label{sub:rose-low}

A dissipative system held far from equilibrium does not generate at a constant rate. It has phases of high throughput and phases of low; it cycles between them, and the low phase is not a malfunction of the system but a part of its cycle. In the low phase the system's visible activity falls---and were one to judge it only by that visible activity, one would read the low phase as failure. But the low phase is frequently the condition of the next high one: the phase in which the gradients that will drive the next surge are re-established, in which what was expended is replenished, in which the system reorganises below the threshold of visible output. A relation has such phases. After a surge of generation---or after the disequilibrium of an estrangement---it may draw inward, slacken its visible coupling, go quiet. This is not, as such, the failure of the bond; it may be the dormancy from which the bond renews.

The error the section warns against is the same one the whole paper has tracked, in its last and gentlest form: the demand that a living system display constant output, that a relation generate visibly and without remission or stand condemned. This demand, applied to the dormant phase, is destructive in a precise way. It treats the gathering of force as a loss of force, and in its anxiety to restore visible activity it disturbs the very quiet in which the gathering occurs---forcing throughput when the system is replenishing, and so preventing the replenishment. To insist that a relation bloom in every season is to forbid it the winter in which the next bloom is prepared. The capacity to let a bond go dormant without reading the dormancy as death is, in its quiet way, among the more difficult of the practical capacities the paper recommends.

\begin{claim}[Dormancy is not decline]
The low, inward phase of a generative relation---the slackening of visible coupling after a surge or an estrangement---is not, as such, the failure of the bond but may be its dormancy: the phase in which expended force is replenished and the gradients of the next surge re-established, below the threshold of visible output. The destructive error is to read the gathering of force as its loss, and, in anxiety to restore visible activity, to force throughput during replenishment and so prevent it. To demand that a relation bloom in every season is to forbid it the winter in which the next bloom is prepared.
\end{claim}

\subsection{The rose in winter}
\label{sub:rose-winter}

The paper began with two bodies falling, and may end the cosmological movement with the other image the series has long carried: the rose. A rose in winter has not died. It has drawn its life down into the root, where nothing flowers and nothing, to the eye, occurs; and the drawing-down is not the opposite of the bloom but its condition---the season in which the rose does the unseen work without which there is no spring. The withering of the visible flower is the storing of force in the unseen root. One who loved the rose only in bloom, and read every winter as a death, would tear up by the roots the very thing whose flowering they wished to keep.

This is offered as an image and not as an argument; the argument is the dissipative one of the preceding paragraphs, and it stands without the figure. But the figure carries something the argument states more coldly: that the phases in which a bond seems to give nothing may be the phases in which it is most quietly faithful to its own renewal, and that the love which can abide the winter of a bond---neither forcing it to bloom nor mistaking its dormancy for its death---is a love that has understood the spiral. The estrangements this paper has anatomised, the passings-without-meeting, the withdrawals and the quiets, are not, in a living bond, the failures of love. They are, at their deepest, love's winters: the inward seasons of a generativity that returns, when it returns, not to where it was but lifted---if the root is left, through the cold, undisturbed.

\noindent
With the cosmological movement closed---disequilibrium as the condition of life, the spiral as the form of return, dormancy as the gathering of force---the paper turns from the widest frame to the narrowest and most practical: what, given all of this, two people are actually to \emph{do}. The praxis that follows is disciplined throughout by the two-level structure these sections have only deepened: that one tends conditions and does not command blooms, and that the work divides into a labour of letting-be and a labour of active defence.
